% !TEX root = ../../main.tex

\section{Evaluation Plan and Missing Results}

This section is a prospective evaluation plan, not a completed results section.
[TODO: replace this section with completed experiments, report sample sizes and
random seeds, and connect each result to the method claims made in the
Introduction and Discussion.]
Existing benchmark scripts and result folders are useful development evidence,
but they are not yet a locked manuscript result. A submission-ready results
section needs a fixed commit, a fixed configuration manifest, random seeds,
case definitions, and output files that can be traced directly from every
reported claim.
An empirical evaluation of Tree-Break Selection should cover four main blocks.

\subsection{Does the Method Stay Calibrated Under the Null?}

The first requirement is a null calibration study on binary data with no true
cluster structure. The goal is to measure how often the child-parent edge test
and the sibling test reject under scenarios where they should not. This should
be reported both at the
individual-test level and at the resulting number-of-clusters level. The latter
is critical because a method can have mild per-test inflation yet still produce
severe over-splitting once the tree walk is applied.

\subsection{Does the Method Recover Planted Structure?}

The second requirement is a controlled power study with planted binary subtrees.
This should vary sample size per cluster, effect size in feature rates, the
number of informative features, tree imbalance, and background noise level.

Power should be measured both for recovery of important tree splits and for the
final clustering outcome. Adjusted Rand Index and exact split recovery are both
useful endpoints.

\subsection{Where Is the Method Robust or Fragile?}

For a paper centered on one implemented method, the key empirical question is
robustness. The most useful stress tests are sensitivity to small child leaf
counts in the child-parent edge test, sparse high-dimensional feature matrices,
tree imbalance, shallow-versus-deep signal placement, and the pass-through rule
used by the final traversal.

These analyses should show where the implemented method is well calibrated,
where it becomes under-powered, and where it tends to over-split or
under-split.

\subsection{How Does the Method Behave on Benchmarks?}

The repository already contains benchmark infrastructure and reporting code. A
manuscript-quality benchmark section should align those experiments with the
method claims in the text, focusing on:

\begin{itemize}[leftmargin=1.5em]
  \item synthetic binary benchmarks with known labels,
  \item sensitivity to over-splitting and under-splitting,
  \item comparison to standard clustering baselines, and
  \item qualitative inspection of annotated trees.
\end{itemize}

\subsection{Which Defaults Need Ablation?}

The defaults in \Cref{tab:method-constants} should not be presented as
validated constants until they are tested. The minimum spectral dimension, the
Marchenko--Pastur dimension rule, the sibling projection-dimension rule, the
context bandwidth, the effective calibration sample size, the edge and sibling
testing levels, and the pass-through traversal should each appear in at least
one ablation or sensitivity analysis.
[VALIDATION GAP: no manuscript-ready ablation table currently validates the
method constants in \Cref{tab:method-constants}.]

\subsection{What Counts as Submission-Ready Evidence?}

A result should enter the paper only when the corresponding script, case
registry, random seeds, software version, and output file are identified. Each
reported number should map to one reproducible result artifact, and each
interpretive claim should cite the table or figure that supports it.
[EMPIRICAL GAP: add a locked evaluation run, a result manifest, and a
manuscript table mapping each claim to the corresponding output file.]
